% Part: first-order-logic
% Chapter: completeness
% Section: lindenbaums-lemma

\documentclass[../../../include/open-logic-section]{subfiles}

\begin{document}

\iftag{FOL}
      {\olfileid{fol}{com}{lin}}
      {\olfileid{pl}{com}{lin}}

\olsection{Lindenbaum引理}

\begin{explain}
我们现在证明一个引理：任何一致的语句集都包含于某个不仅一致而且完备的语句集中。
证明方法是每次添加一个语句，并保证在每一步集合都保持一致。
这样做的目的是使得对于每个 $!A$，总有 $!A$ 或 $\lnot !A$ 在某个阶段被添加进去。
该构造中所有阶段的并集将包含 $!A$ 或其否定 $\lnot !A$，因而是完备的。
它也是一致的，因为我们确保了在每个阶段都不引入不一致。
\end{explain}

\begin{lem}[Lindenbaum 引理]
\ollabel{lem:lindenbaum} 语言~$\Lang{L}$ 中的每个一致集合~$\Gamma$ 都可以扩张为一个完备且一致的集合~$\Gamma^*$。
\end{lem}

\begin{proof}
设 $\Gamma$ 是一致的。令 $!A_0$, $!A_1$, \dots{} 为 $\Lang L$ 中所有语句的一个枚举。
定义 $\Gamma_0 = \Gamma$，并 \[
\Gamma_{n+1} =
\begin{cases}
\Gamma_n \cup \{ !A_n \} & \textrm{if $\Gamma_n \cup \{!A_n\}$ is
  consistent;} \\
\Gamma_n \cup \{ \lnot !A_n \} & \textrm{otherwise.}
\end{cases}
\] 令 $\Gamma^* = \bigcup_{n \geq 0} \Gamma_n$。

每个 $\Gamma_n$ 都是一致的：$\Gamma_0$ 根据定义是一致的。
若 $\Gamma_{n+1} = \Gamma_n \cup \{!A_n\}$，这是因为后者是一致的；否则，$\Gamma_{n+1} = \Gamma_n \cup \{\lnot
!A_n\}$。
我们需要验证 $\Gamma_n \cup \{\lnot !A_n\}$ 是一致的。
假设并非如此。那么 $\Gamma_n \cup
\{!A_n\}$ \emph{和} $\Gamma_n \cup \{\lnot !A_n\}$ 都是不一致的。
这意味着由
\iftag{FOL}{%
  \tagrefs{prfAX/{fol:axd:prv:prop:provability-exhaustive},
    prfSC/{fol:seq:prv:prop:provability-exhaustive},
    prfND/{fol:ntd:prv:prop:provability-exhaustive},
    prfTab/{fol:tab:prv:prop:provability-exhaustive}}}{%
  \tagrefs{prfAX/{pl:axd:prv:prop:provability-exhaustive},
    prfSC/{pl:seq:prv:prop:provability-exhaustive},
    prfND/{pl:ntd:prv:prop:provability-exhaustive},
    prfTab/{pl:tab:prv:prop:provability-exhaustive}}}
可得 $\Gamma_n$ 将是不一致的，这与归纳假设矛盾。

对于每个 $n$ 和每个 $i < n$，$\Gamma_i \subseteq \Gamma_n$。
这通过对 $n$ 的简单归纳可得。对于 $n=0$，不存在 $i < 0$，故该断言自动成立。
对于归纳步骤，假设它对 $n$ 成立。我们证明若 $i < n+1$，则 $\Gamma_i \subseteq
\Gamma_{n+1}$。
根据构造，我们有 $\Gamma_{n+1} = \Gamma_n \cup \{!A_n\}$ 或 $=
\Gamma_n \cup \{\lnot !A_n\}$。因此 $\Gamma_n \subseteq
\Gamma_{n+1}$。
若 $i < n+1$，则由归纳假设（若 $i < n$）或 $\Gamma_n \subseteq \Gamma_n$ 这一平凡事实（若 $i = n$），可得 $\Gamma_i \subseteq \Gamma_n$。
由 $\subseteq$ 的传递性，我们得到 $\Gamma_i \subseteq
\Gamma_{n+1}$。

由此可推出 $\Gamma^*$ 是一致的。理由如下：设 $\Gamma' \subseteq \Gamma^*$ 是有限的。
每个 $!B \in \Gamma'$ 也都在某个 $\Gamma_i$ 中。令 $n$ 为其中最大的 $i$。
因为若 $i \le n$ 则 $\Gamma_i \subseteq \Gamma_n$，所以每个 $!B \in \Gamma'$ 也在 $\Gamma_n$ 中，
即 $\Gamma' \subseteq \Gamma_n$，且 $\Gamma_n$ 是一致的。
因此，每个有限子集 $\Gamma' \subseteq \Gamma^*$ 都是一致的。
由 \iftag{FOL}{%
  \tagrefs{prfAX/{fol:axd:ptn:prop:proves-compact},
    prfSC/{fol:seq:ptn:prop:proves-compact},
    prfND/{fol:ntd:ptn:prop:proves-compact},
    prfTab/{fol:tab:ptn:prop:proves-compact}}}{%
  \tagrefs{prfAX/{pl:axd:ptn:prop:proves-compact},
    prfSC/{pl:seq:ptn:prop:proves-compact},
    prfND/{pl:ntd:ptn:prop:proves-compact},
    prfTab/{pl:tab:ptn:prop:proves-compact}}}，$\Gamma^*$ 是一致的。

$\Frm[L]$ 中的每个语句都出现在用于定义 $\Gamma^*$ 的列表中。
如果 $!A_n \notin \Gamma^*$，则是因为 $\Gamma_n \cup \{!A_n\}$ 是不一致的。
但由此可知 $\lnot !A_n \in \Gamma^*$，所以 $\Gamma^*$ 是完备的。
\end{proof}

\end{document}